% =========================================================
% Monotone Sequences — Curated Exercise Set
% Sources: Abbott, Bartle & Sherbert, Bruckner, Tao, Lebl
% =========================================================

\section*{Monotone Sequences — Exercise Review Set}

Monotone sequences provide the first major application of the completeness
property of the real numbers. These exercises train the techniques used in
monotone convergence arguments.

\bigskip

% ---------------------------------------------------------
% Stub Macro
% ---------------------------------------------------------

\newcommand{\exstub}[4]{
\subsection*{Exercise: #1}

\textbf{Source:} #2

\textbf{Technique:} #3

\textbf{Task.} #4

\begin{remark*}[Work Area]
Write your proof or computation here.
\end{remark*}

\bigskip
}

% =========================================================
% Abbott
% =========================================================

\section*{Exercises from Abbott}

\exstub
{Abbott 2.3.8}
{Abbott, \textit{Understanding Analysis}}
{monotone-bound argument}
{Prove that every bounded monotone sequence converges.}

\exstub
{Abbott 2.3.9}
{Abbott}
{supremum construction}
{Let $(a_n)$ be increasing and bounded above.
Show that $\lim a_n = \sup\{a_n\}$.}

\exstub
{Abbott 2.3.10}
{Abbott}
{monotone-bound argument}
{Give an example of a monotone sequence that diverges.}

\exstub
{Abbott 2.3.11}
{Abbott}
{sequence bounds}
{Show that every increasing sequence bounded above
has a finite limit.}

\exstub
{Abbott 2.3.12}
{Abbott}
{epsilon-approximation}
{Let $s=\sup\{a_n\}$. Show that for every $\varepsilon>0$
there exists $N$ such that $s-\varepsilon < a_n \le s$ for $n\ge N$.}

% =========================================================
% Bartle
% =========================================================

\section*{Exercises from Bartle \& Sherbert}

\exstub
{Bartle 3.4.1}
{Bartle \& Sherbert, \textit{Introduction to Real Analysis}}
{monotone-bound argument}
{Prove that every bounded increasing sequence converges.}

\exstub
{Bartle 3.4.2}
{Bartle \& Sherbert}
{supremum construction}
{Show that the limit of a bounded increasing sequence
equals the supremum of its terms.}

\exstub
{Bartle 3.4.3}
{Bartle \& Sherbert}
{infimum construction}
{Show that the limit of a bounded decreasing sequence
equals the infimum of its terms.}

\exstub
{Bartle 3.4.5}
{Bartle \& Sherbert}
{monotone-bound argument}
{Show that if a sequence is monotone and bounded,
then it converges.}

\exstub
{Bartle 3.4.7}
{Bartle \& Sherbert}
{tail argument}
{If $(a_n)$ is increasing and bounded,
prove $(a_{n+k})$ has the same limit.}

% =========================================================
% Bruckner
% =========================================================

\section*{Exercises from Bruckner}

\exstub
{Bruckner 2.2.1}
{Bruckner, \textit{Real Analysis}}
{monotone-bound argument}
{Prove the monotone convergence theorem for sequences.}

\exstub
{Bruckner 2.2.2}
{Bruckner}
{supremum construction}
{Let $(a_n)$ be increasing.
Show that if it is bounded above,
then $\lim a_n = \sup\{a_n\}$.}

\exstub
{Bruckner 2.2.3}
{Bruckner}
{infimum construction}
{Prove the analogous result for decreasing sequences.}

\exstub
{Bruckner 2.2.4}
{Bruckner}
{squeeze argument}
{Let $(a_n)$ and $(b_n)$ be monotone sequences with
$a_n \le c_n \le b_n$. Show that if
$a_n \to L$ and $b_n \to L$ then $c_n \to L$.}

\exstub
{Bruckner 2.2.6}
{Bruckner}
{sequence bounds}
{Give an example of a monotone unbounded sequence.}

% =========================================================
% Tao
% =========================================================

\section*{Exercises from Tao}

\exstub
{Tao 6.3.6}
{Tao, \textit{Analysis I}}
{monotone-bound argument}
{Prove the monotone convergence theorem.}

\exstub
{Tao 6.3.7}
{Tao}
{supremum construction}
{Let $(a_n)$ be increasing.
Show $\lim a_n = \sup\{a_n\}$.}

\exstub
{Tao 6.3.8}
{Tao}
{sequence bounds}
{Give an example of a monotone sequence
whose limit is infinite.}

\exstub
{Tao 6.3.9}
{Tao}
{tail argument}
{Show that if a monotone sequence converges,
then every subsequence converges to the same limit.}

\exstub
{Tao 6.3.10}
{Tao}
{subsequence extraction}
{Show that every bounded sequence has
a monotone subsequence.}

% =========================================================
% Lebl
% =========================================================

\section*{Exercises from Lebl}

\exstub
{Lebl 2.2.4}
{Lebl, \textit{Basic Analysis I}}
{monotone-bound argument}
{Show that every bounded monotone sequence converges.}

\exstub
{Lebl 2.2.5}
{Lebl}
{supremum construction}
{Let $(a_n)$ be increasing and bounded above.
Show $\lim a_n = \sup\{a_n\}$.}

\exstub
{Lebl 2.2.6}
{Lebl}
{infimum construction}
{Prove the analogous result for decreasing sequences.}

\exstub
{Lebl 2.2.7}
{Lebl}
{squeeze argument}
{Let $a_n \le b_n \le c_n$.
If $a_n$ and $c_n$ converge to the same limit,
show $b_n$ converges.}

\exstub
{Lebl 2.2.9}
{Lebl}
{sequence bounds}
{Show that every monotone sequence has a limit
or diverges to $\pm\infty$.}

